%==============================================================================
% sections/00_abstract.tex
%==============================================================================
\begin{abstract}
\noindent
We study rare-event probabilities and efficient simulation for finite sums of
regularly varying factor-loss contributions. The first layer of the paper gives
sharp independent-factor theory. For independent, non-identically distributed
summands satisfying pointwise tail equivalence, we prove the exact finite-
portfolio catastrophe principle for the maximum and the corresponding single-
big-jump equivalence for the sum under the stated truncation and left-tail
conditions. Under direct second-order tail equivalence, local-shift regularity,
and finite first moments, we derive a corrected expansion for
$\Pb(S_N>x)$. The mean-shift correction is a leave-one-out term,
$\mu_{-i}=\sum_{j\ne i}\E X_j$, because the large summand supplies the
threshold crossing while the remaining portfolio supplies the first-order
translation. This form passes the $N=1$ consistency check and produces explicit
Pareto/Lomax benchmarks.

The second layer develops a staged extension path for non-independent factors.
For arbitrary dependence, we prove an exact disjoint-partition conditional Monte
Carlo identity in which the independent marginal survival function is replaced
by the conditional tail $\Pb(X_i>T_i\mid X_{-i})$. We then specialize this
identity to four dependence structures that cover the main financial use cases:
latent independent shocks and common shocks, independent blocks with dependent
within-block components, copula and vine-copula conditional tails, and
multivariate regular variation with a spectral/radial representation. The
resulting hierarchy replaces the independent ``one large coordinate'' principle
by ``one large latent shock'', ``one large block'', and ultimately ``one large
spectral vector shock''. Hidden regular variation is formulated as a second-
order joint-tail correction when ordinary multivariate regular variation puts
first-order mass on the coordinate axes but slower joint-tail cones remain
empirically material.

Building on these asymptotics, we analyze summed CdMC, single-index stratified
CdMC, Pareto/Lomax closed forms, and centered control-variate estimators. In
the independent case the summed CdMC estimator is unbiased and has the scale-
invariant bounded-relative-error bound $N^\alpha-1$. In dependent settings the
same partition gives unbiasedness, while bounded-relative-error guarantees
require model-specific conditional-tail, block, latent-shock, or radial-envelope
assumptions. The manuscript includes full proof templates for the dependent
identities and for the latent-shock, block, copula, and spectral/radial
extensions.

The numerical and empirical sections are written as a reproducible long-form
analysis backed by a complete reference implementation. The simulation suite
exercises six data-generating families (independent inid, common-shock,
block, copula, MRV with radial-angular separation, and hidden regular
variation) and is summarised in
\cref{fig:ind-tail-equivalence,fig:max-vs-sum,fig:exp-eff-vs-x,fig:exp-eff-vs-alpha,fig:exp-eff-vs-amin,fig:relative-error,fig:stratified,fig:second-order-placeholder,fig:common-shock-geometry,fig:spectral-simplex-placeholder,fig:hidden-cones-placeholder,fig:simulation-dashboard-placeholder,fig:bootstrap-audit}.
The headline simulation findings are: (i) the corrected second-order
expansion reduces absolute relative error by roughly an order of magnitude
over the leading sum-tail asymptotic (\cref{fig:second-order-placeholder});
(ii) for heterogeneous-$\alpha$ designs the BRE-bounded efficiency rate is
controlled by $\alpha_{\min}$ rather than $\bar\alpha$
(\cref{fig:exp-eff-vs-alpha,fig:exp-eff-vs-amin}); (iii) the choice of
control-variate surrogate dominates the choice of pilot rule by orders of
magnitude in $\widehat\rho^2$
(\cref{fig:relative-error}); (iv) iid bootstrap bands for the spectral
constant under-cover by 60 percentage points when the series carries
moderate autocorrelation, while block and stationary bootstraps remain at
nominal coverage (\cref{fig:bootstrap-audit}).

The real-data design uses the public Kenneth-French Data Library
(1926-07-01 to 2026-03-31, $n = 26\,212$ daily observations). Hill / POT
tail-index estimators and the pairwise $\chi$, $\bar\chi$, $\eta$
dependence diagnostics produce numbers consistent with the published
literature on US equity factor returns
\citep{Gabaix2009,Cont2001,EmbrechtsMcNeilFrey2005,ChristoffersenErrunzaJacobsLangloi2012}.
A separate companion document
(\texttt{docs/REAL\_DATA\_VALIDATION.md} in the
\texttt{FactorTail} repository) records the exact numbers and the
literature anchors. The crisis-vs-full sample comparison shows that long
panels (1926--2026) include events more extreme than the 2008 financial
crisis, yielding Hill estimates slightly below Gabaix's published $\sim 3$.

\medskip\noindent The paper positions itself as a sharpening and extension of
the independent-factor regularly-varying rare-event framework of
\citet{PourbabaeeSolari2019}, integrating tools from heavy-tailed rare-event
simulation \citep{AsmussenKroese2006,HartingerKortschak2009}, multivariate
extremes \citep{HultLindskogMikoschSamorodnitsky2005,Resnick2007}, hidden
regular variation \citep{LedfordTawn1996,MaulikResnick2005}, and dependent
rare-event simulation
\citep{BlanchetRojasNandayapa2011,ChengFuhPang2025} into a single staged
estimator-and-diagnostic pipeline. An annotated literature synthesis appears
in \cref{sec:literature}. The methods are implemented in the open-source
\texttt{FactorTail} library
(\url{https://github.com/osolari/FactorTail}), which provides reference
implementations of every estimator, diagnostic, and reproducibility artifact
discussed in this paper.
\end{abstract}
